\documentclass[11pt]{article}
\newcommand{\ListaTitulo}{Lista 13 --- Fatoriais $2^k$}
% Preâmbulo compartilhado — MATD48, Listas de Exercícios 2026
% Usado por Lista{NN}.tex e Gabarito{NN}.tex via % Preâmbulo compartilhado — MATD48, Listas de Exercícios 2026
% Usado por Lista{NN}.tex e Gabarito{NN}.tex via % Preâmbulo compartilhado — MATD48, Listas de Exercícios 2026
% Usado por Lista{NN}.tex e Gabarito{NN}.tex via \input{preamble}

\usepackage[utf8]{inputenc}
\usepackage[T1]{fontenc}
\usepackage[brazil]{babel}
\usepackage{fullpage}
\usepackage{amsmath,amssymb}
\usepackage{enumitem}
\usepackage{xcolor}
\usepackage{fancyhdr}
\usepackage{titlesec}
\usepackage{listings}
\usepackage{hyperref}
\usepackage{booktabs}
\usepackage{graphicx}

\definecolor{matd48azul}{HTML}{0B4F6C}
\definecolor{matd48verde}{HTML}{2E8B57}

\titleformat{\section}{\normalfont\Large\bfseries\color{matd48azul}}{\thesection}{1em}{}
\titleformat{\subsection}{\normalfont\large\bfseries\color{matd48azul}}{\thesubsection}{1em}{}

\providecommand{\ListaTitulo}{Lista de Exercícios}

\setlength{\headheight}{14pt}
\pagestyle{fancy}
\fancyhf{}
\lhead{\small MATD48 --- Planejamento de Experimentos A}
\rhead{\small \ListaTitulo}
\cfoot{\thepage}
\renewcommand{\headrulewidth}{0.4pt}

\lstset{
  basicstyle=\ttfamily\small,
  breaklines=true,
  frame=single,
  columns=fullflexible,
  backgroundcolor=\color{gray!8},
  commentstyle=\color{matd48verde},
  keywordstyle=\color{matd48azul}\bfseries,
  language=R,
  extendedchars=true,
  literate=%
    {á}{{\'a}}1 {à}{{\`a}}1 {â}{{\^a}}1 {ã}{{\~a}}1
    {é}{{\'e}}1 {ê}{{\^e}}1 {í}{{\'i}}1 {ó}{{\'o}}1
    {ô}{{\^o}}1 {õ}{{\~o}}1 {ú}{{\'u}}1 {ç}{{\c c}}1
    {Á}{{\'A}}1 {À}{{\`A}}1 {Â}{{\^A}}1 {Ã}{{\~A}}1
    {É}{{\'E}}1 {Ê}{{\^E}}1 {Í}{{\'I}}1 {Ó}{{\'O}}1
    {Ô}{{\^O}}1 {Õ}{{\~O}}1 {Ú}{{\'U}}1 {Ç}{{\c C}}1
}

\hypersetup{colorlinks=true, linkcolor=matd48azul, urlcolor=matd48azul, citecolor=matd48azul}

\newcommand{\cabecalho}[2]{%
  \begin{center}
    {\Large\bfseries\color{matd48azul} #1}\\[0.3em]
    {\large #2}\\[0.3em]
    \rule{\linewidth}{0.6pt}
  \end{center}
  \vspace{0.5em}
}

\newenvironment{questao}[1]{\vspace{1em}\noindent\textbf{Questão #1.}\hspace{0.4em}}{\par}
\newenvironment{solucao}{\vspace{0.3em}\noindent\textit{Solução.}\hspace{0.4em}\color{black!85}}{\par\vspace{0.5em}}


\usepackage[utf8]{inputenc}
\usepackage[T1]{fontenc}
\usepackage[brazil]{babel}
\usepackage{fullpage}
\usepackage{amsmath,amssymb}
\usepackage{enumitem}
\usepackage{xcolor}
\usepackage{fancyhdr}
\usepackage{titlesec}
\usepackage{listings}
\usepackage{hyperref}
\usepackage{booktabs}
\usepackage{graphicx}

\definecolor{matd48azul}{HTML}{0B4F6C}
\definecolor{matd48verde}{HTML}{2E8B57}

\titleformat{\section}{\normalfont\Large\bfseries\color{matd48azul}}{\thesection}{1em}{}
\titleformat{\subsection}{\normalfont\large\bfseries\color{matd48azul}}{\thesubsection}{1em}{}

\providecommand{\ListaTitulo}{Lista de Exercícios}

\setlength{\headheight}{14pt}
\pagestyle{fancy}
\fancyhf{}
\lhead{\small MATD48 --- Planejamento de Experimentos A}
\rhead{\small \ListaTitulo}
\cfoot{\thepage}
\renewcommand{\headrulewidth}{0.4pt}

\lstset{
  basicstyle=\ttfamily\small,
  breaklines=true,
  frame=single,
  columns=fullflexible,
  backgroundcolor=\color{gray!8},
  commentstyle=\color{matd48verde},
  keywordstyle=\color{matd48azul}\bfseries,
  language=R,
  extendedchars=true,
  literate=%
    {á}{{\'a}}1 {à}{{\`a}}1 {â}{{\^a}}1 {ã}{{\~a}}1
    {é}{{\'e}}1 {ê}{{\^e}}1 {í}{{\'i}}1 {ó}{{\'o}}1
    {ô}{{\^o}}1 {õ}{{\~o}}1 {ú}{{\'u}}1 {ç}{{\c c}}1
    {Á}{{\'A}}1 {À}{{\`A}}1 {Â}{{\^A}}1 {Ã}{{\~A}}1
    {É}{{\'E}}1 {Ê}{{\^E}}1 {Í}{{\'I}}1 {Ó}{{\'O}}1
    {Ô}{{\^O}}1 {Õ}{{\~O}}1 {Ú}{{\'U}}1 {Ç}{{\c C}}1
}

\hypersetup{colorlinks=true, linkcolor=matd48azul, urlcolor=matd48azul, citecolor=matd48azul}

\newcommand{\cabecalho}[2]{%
  \begin{center}
    {\Large\bfseries\color{matd48azul} #1}\\[0.3em]
    {\large #2}\\[0.3em]
    \rule{\linewidth}{0.6pt}
  \end{center}
  \vspace{0.5em}
}

\newenvironment{questao}[1]{\vspace{1em}\noindent\textbf{Questão #1.}\hspace{0.4em}}{\par}
\newenvironment{solucao}{\vspace{0.3em}\noindent\textit{Solução.}\hspace{0.4em}\color{black!85}}{\par\vspace{0.5em}}


\usepackage[utf8]{inputenc}
\usepackage[T1]{fontenc}
\usepackage[brazil]{babel}
\usepackage{fullpage}
\usepackage{amsmath,amssymb}
\usepackage{enumitem}
\usepackage{xcolor}
\usepackage{fancyhdr}
\usepackage{titlesec}
\usepackage{listings}
\usepackage{hyperref}
\usepackage{booktabs}
\usepackage{graphicx}

\definecolor{matd48azul}{HTML}{0B4F6C}
\definecolor{matd48verde}{HTML}{2E8B57}

\titleformat{\section}{\normalfont\Large\bfseries\color{matd48azul}}{\thesection}{1em}{}
\titleformat{\subsection}{\normalfont\large\bfseries\color{matd48azul}}{\thesubsection}{1em}{}

\providecommand{\ListaTitulo}{Lista de Exercícios}

\setlength{\headheight}{14pt}
\pagestyle{fancy}
\fancyhf{}
\lhead{\small MATD48 --- Planejamento de Experimentos A}
\rhead{\small \ListaTitulo}
\cfoot{\thepage}
\renewcommand{\headrulewidth}{0.4pt}

\lstset{
  basicstyle=\ttfamily\small,
  breaklines=true,
  frame=single,
  columns=fullflexible,
  backgroundcolor=\color{gray!8},
  commentstyle=\color{matd48verde},
  keywordstyle=\color{matd48azul}\bfseries,
  language=R,
  extendedchars=true,
  literate=%
    {á}{{\'a}}1 {à}{{\`a}}1 {â}{{\^a}}1 {ã}{{\~a}}1
    {é}{{\'e}}1 {ê}{{\^e}}1 {í}{{\'i}}1 {ó}{{\'o}}1
    {ô}{{\^o}}1 {õ}{{\~o}}1 {ú}{{\'u}}1 {ç}{{\c c}}1
    {Á}{{\'A}}1 {À}{{\`A}}1 {Â}{{\^A}}1 {Ã}{{\~A}}1
    {É}{{\'E}}1 {Ê}{{\^E}}1 {Í}{{\'I}}1 {Ó}{{\'O}}1
    {Ô}{{\^O}}1 {Õ}{{\~O}}1 {Ú}{{\'U}}1 {Ç}{{\c C}}1
}

\hypersetup{colorlinks=true, linkcolor=matd48azul, urlcolor=matd48azul, citecolor=matd48azul}

\newcommand{\cabecalho}[2]{%
  \begin{center}
    {\Large\bfseries\color{matd48azul} #1}\\[0.3em]
    {\large #2}\\[0.3em]
    \rule{\linewidth}{0.6pt}
  \end{center}
  \vspace{0.5em}
}

\newenvironment{questao}[1]{\vspace{1em}\noindent\textbf{Questão #1.}\hspace{0.4em}}{\par}
\newenvironment{solucao}{\vspace{0.3em}\noindent\textit{Solução.}\hspace{0.4em}\color{black!85}}{\par\vspace{0.5em}}


\begin{document}

\cabecalho{Lista de Exercícios 13}{Fatoriais $2^k$: ortogonalidade, análise sem repetição e confusão (Capítulo 6 / Aula 13)}

\noindent \textit{Entrega: até a próxima aula. Justifique todas as respostas --- nestas listas, a
justificativa vale mais que a resposta final. O gabarito completo está em \texttt{Gabarito13.pdf}.}

\begin{questao}{1} \textbf{(Dedução --- ortogonalidade de um $2^3$ e o estimador de efeito)}

Considere um fatorial $2^3$ completo, uma repetição ($n=2^3=8$ corridas), fatores codificados
$\pm1$. Seja $\mathbf{x}_j \in \{-1,+1\}^8$ a coluna de sinais de um efeito qualquer (principal ou
de interação, construída por produto de colunas, Seção 6.1 do livro-texto), e
$\mathbf{X} = [\mathbf{1}\ |\ \mathbf{x}_A\ |\ \mathbf{x}_B\ |\ \cdots\ |\ \mathbf{x}_{ABC}]$ a
matriz de delineamento completa (8 colunas).

\begin{enumerate}[label=(\alph*)]
  \item Liste as 8 combinações de sinais $(x_A,x_B,x_C)$ e mostre, por soma direta, que
    $\mathbf{x}_A'\mathbf{x}_B = 0$.
  \item Mostre que $\mathbf{x}_A'\mathbf{x}_{BC} = 0$, onde $x_{BC}=x_Bx_C$. (Dica: escreva o
    produto $x_A\cdot x_Bx_C$ e observe que ele percorre os valores $\pm1$ em igual proporção.)
  \item Argumente, generalizando os itens (a)-(b), por que \emph{toda} matriz $\mathbf{X}$ de um
    $2^k$ completo com uma repetição, codificada em $\pm1$, satisfaz $\mathbf{X}'\mathbf{X} =
    n\,\mathbf{I}_{2^k}$ ($n\,$ vezes a matriz identidade).
  \item A partir de $\mathbf{X}'\mathbf{X}=n\mathbf{I}$, mostre que $\hat\beta_j =
    \mathbf{x}_j'\mathbf{Y}/n$ e que, portanto, o efeito estimado é
    $$\text{efeito}_j = 2\hat\beta_j = 2^{-(k-1)}\sum_{i=1}^{2^k} (\pm1)\,y_i,$$
    a fórmula do contraste de Yates. (Dica: $2/n = 2/2^k = 2^{1-k} = 2^{-(k-1)}$.)
\end{enumerate}
\end{questao}

\begin{questao}{2} \textbf{(Ciência de dados --- calculando efeitos pela tabela de sinais)}

Um experimento de engenharia mede o tempo de build (min) de um pipeline de CI/CD sob três fatores
de configuração, cada um em dois níveis: $A$ = cache de dependências (desligado/ligado), $B$ =
paralelismo de testes (baixo/alto), $C$ = otimização de compilação (desligada/ligada). O fatorial
$2^3$ completo (uma corrida por combinação) resultou em:

\begin{center}
\begin{tabular}{c|cccccccc}
Corrida & $(1)$ & $a$ & $b$ & $ab$ & $c$ & $ac$ & $bc$ & $abc$ \\ \hline
Tempo (min) & 20 & 24 & 28 & 36 & 16 & 20 & 24 & 32
\end{tabular}
\end{center}

\begin{enumerate}[label=(\alph*)]
  \item Construa a tabela de sinais completa ($A$, $B$, $C$, $AB$, $AC$, $BC$, $ABC$) para as 8
    corridas, na ordem padrão de Yates dada acima.
  \item Usando a fórmula do contraste (Questão 1d) com $k=3$ (logo $2^{-(k-1)}=1/4$), calcule os
    7 efeitos: $A$, $B$, $C$, $AB$, $AC$, $BC$, $ABC$.
  \item Sem fazer nenhuma conta nova, ordene os 7 efeitos por magnitude absoluta, do maior para o
    menor.
\end{enumerate}
\end{questao}

\begin{questao}{3} \textbf{(Ciência de dados --- método de Lenth na mão)}

Usando os 7 efeitos que você calculou na Questão 2 (não há repetição — este pipeline de CI/CD
só pôde ser executado uma vez por configuração, por custo de infraestrutura):

\begin{enumerate}[label=(\alph*)]
  \item Calcule $s_0 = 1{,}5\cdot\text{mediana}_i|c_i|$.
  \item Calcule o PSE $=1{,}5\cdot\text{mediana}$ dos $|c_i|$ que satisfazem $|c_i| < 2{,}5\,s_0$.
  \item Com $m=7$ efeitos, $d=m/3$. Usando $t_{0{,}975;\,2{,}33}\approx3{,}76$ (fornecido), calcule
    a margem de erro $\text{ME} = t_{0{,}975;d}\cdot\text{PSE}$.
  \item Quais efeitos são declarados "grandes" (i.e., $|c_i|>\text{ME}$)? Isso é consistente com o
    que você esperaria observando a tabela de efeitos da Questão 2(c)?
\end{enumerate}
\end{questao}

\begin{questao}{4} \textbf{(Confusão --- designação de blocos e colinearidade)}

O pipeline de CI/CD da Questão 2 precisa, na prática, ser executado em dois lotes de máquinas
virtuais diferentes (um "bloco" de infraestrutura por lote), 4 corridas por lote.

\begin{enumerate}[label=(\alph*)]
  \item Sabendo que o efeito $ABC$ foi estimado como (praticamente) nulo na Questão 2, proponha um
    esquema de confusão que sacrifique $ABC$: calcule o sinal de $ABC$ para cada uma das 8
    corridas e divida-as em dois blocos de 4.
  \item Em termos da matriz $\mathbf{X}$, explique por que a coluna de bloco que você construiu no
    item (a) fica \emph{idêntica} (a menos de sinal) à coluna $\mathbf{x}_{ABC}$, e por que isso
    torna impossível separar "efeito de bloco" de "interação $ABC$" a partir dos dados.
  \item Se fosse possível rodar \emph{duas} repetições completas do $2^3$ (16 corridas, em 4
    lotes de 4), por que usar geradores diferentes em cada repetição --- por exemplo, $ABC$ na
    repetição 1 e $BC$ na repetição 2 --- é preferível a repetir o mesmo gerador $ABC$ nas duas?
\end{enumerate}
\end{questao}

\begin{questao}{5} \textbf{(Ciência de dados --- exercício computacional em R)}

O código abaixo automatiza o cálculo de efeitos e do método de Lenth para os dados da Questão 2.
Está incompleto.

\begin{lstlisting}
dados <- expand.grid(C = c(-1,1), B = c(-1,1), A = c(-1,1))
dados$y <- c(20, 16, 28, 24, 24, 20, 36, 32)   # reordene se necessario

# (a) ajuste o modelo saturado e extraia os efeitos (2 * coeficientes, sem o intercepto)
modelo <- lm(______________________, data = dados)
efeitos <- 2 * coef(modelo)[-1]

# (b) complete o metodo de Lenth
s0  <- 1.5 * median(abs(efeitos))
pse <- 1.5 * median(______________________)
me  <- qt(0.975, length(efeitos)/3) * pse

# (c) quais efeitos ultrapassam a margem de erro?
efeitos[______________________]
\end{lstlisting}

\begin{enumerate}[label=(\alph*)]
  \item Complete as três linhas indicadas.
  \item Rodando o código, os valores de \texttt{efeitos} devem coincidir com os da Questão 2(b), e
    o resultado do item (c) deve coincidir com sua resposta na Questão 3(d). Aponte uma fonte
    possível de pequena discrepância numérica entre o \texttt{t} exato calculado por
    \texttt{qt()} e o valor arredondado $3{,}76$ fornecido na Questão 3(c).
\end{enumerate}
\end{questao}

\end{document}
